\section{Adornment of the Spiritual Marriage}

A careful reading of \textbf{John of Ruysbroeck}`s works reveals the great distance that separates contemporary thought from Medieval spirituality. The meaning of those older texts is often obscured by translations which assume that archaic language is more uplifting than clear writing. For example, here is a sentence as it appears in the English translation of John's \emph{Adornment of the Spiritual Marriage}\footnote{\url{https://www.ccel.org/ccel/ruysbroeck/adornment.html}}:

\begin{quotex}
[Christ] fed in ghostly wise with true and inward teachings all those men who could understand them: and others from without through the senses with signs and wonders. 

\end{quotex}
Here is my alternate translation for your consideration:

\begin{quotex}
Christ fed the true esoteric teachings in a spiritual way to all those men who could understand them; and others exoterically through the senses with signs and wonders. 

\end{quotex}
Here John is making clear that the esoteric (or inward) teachings have to do with changes in man's interior life, not with unusual events in the space-time world. More specifically, the esoteric teachings require a change in the way one lives his life. A fortiori, these teachings are not part of some novel philosophical system. In the meditation on Justice, \textbf{Valentin Tomberg} brings up the same issues. The spiritual ``Greeks" are interested in worldly wisdom and spiritual ``Jews" expect the power of God to be revealed in signs and wonders. (Obviously, this has nothing to do with ethnicity.)

Another factor to consider is the difference between literal and symbolic interpretations of sacred texts. The exoterically minded will understand symbolic teachings literally and literal teachings symbolically. To illustrate this, I will mention two topics brought up by John and the animal powers and his understanding of the three comings of Christ.

\paragraph{Animal Powers}
In Chapter XXIII, John lists three adversaries to the spiritual life:

\begin{itemize}
\item \textbf{The Devil}. These are spiritual influences that impinge on us from outside the world. 
\item \textbf{The World}. These are the hodgepodge of influences originating in the human world. These are the ``A" influences described by \textbf{Boris Mouravieff}. 
\item \textbf{Our own Flesh}. These are the influences arising from the lower centers of the soul. 
\end{itemize}
For this illustration, we will focus on the ``flesh". This sentence summarizes John's thought on this matter:

\begin{quotex}
A man should keep his senses in sobriety and should restrain the animal powers by means of the reason; so that the lusts of the flesh do not enter too far into the savouring of food and of drink. 

\end{quotex}
Of course, elsewhere he includes sexual impulses along with food and drink. Now this is all to be understood quite literally. ``Animal powers" is simply the admission that man's soul includes an ``animal soul". Moreover, the traditional teaching is that the ``reason", or intellectual soul, is to have dominance over the lower parts of the soul. This is no different from what Rene Guenon or Julius Evola have written about many times, in many ways. The only difference is that John provides us with the specifics about how that is to be accomplished.

Hence, there is nothing inherently evil about the ``flesh" or ``animal powers". Man is part of organic life on earth, and hence falls under the General Law of Life: motivation by fear, sex, and hunger. This was common knowledge at one time; even the Tantric teachings understood them to be embodied in the three lower chakras.

So the problem arises when animal powers become compulsive, and not subject to the dominance of reason. This is quite opposed to the modern view, which understands the power of reason as a late addition to the genetic makeup of man. In this view, the power of reason is the passive force, in service to the needs of sexual and alimentary desires. In this view, therefore, such a use of reason makes a species more ``fit" for survival. Moreover, many believe today that their sexuality represents the deepest part of themselves, so it needs to be expressed without any \emph{frein vital}, or inner check.

Needless to say, the esoteric teaching is that the intellect needs to be the active force. The intellect, therefore, needs to restrain the animal powers. The intuitive intelligence, the knowledge of the heart, is the deepest part of man, although it usually remains hidden and merely virtual. The animal soul, then, is just the outer husk.

Now, the animal powers already have an intelligence of their own. Even the material world has an intelligence, otherwise the idea that matter ``obeys" certain laws would make no sense.

In the animal world, a natural intelligence regulates sex and hunger, since they, too, are created by the Logos. Thus most animals have a period of estrus during which mating occurs. Often, such mating is very ritualized and ruled by instinct. The same applies to food. Most animals have specific and limited diets which suit them and are instinctually known.

For man, the situation is different. The adult female is sexually receptive at all times and not just during certain times of the year. Moreover, man is omnivorous, able to eat and digest a wide variety of foods. This suggests that man's ``animal powers" are not intended to function by a separate intelligence as is the case with animals, but rather by man's unified reason. This is not the place to go into all the details, which can be summarized by this note from John:

\begin{quotex}
We should build up a wall and make a separation within ourselves. And the lower part of ourselves, which is beastly and contrary to the virtues, and which wills our separation from God, we should hate and persecute, and we should torment it by means of penances and austerity of life; so that it be always repressed, and subject to reason, that thereby righteousness and purity of heart may always have the upper hand in all the works of virtue. 

\end{quotex}
The wall sounds like the ``cage" described by Mouravieff.

\paragraph{Freedom from Images}
Clearly, sex and food are not bad, since they are necessary for life. Nevertheless, the lower part has no interest in spiritual pursuits, so it resists the dominance of the intellectual soul. This prevents the unity of the soul and turns a man into a mass of competing interests. John relies on the grace of God:

\begin{quotex}
There springs from the side of man … a gathering together of all inward and outward powers in the unity of the spirit, in the bonds of love. 

\end{quotex}
The animal powers hinder this unity through the power of sensible images and fantasies. In other words, images will arise in consciousness of sex and food, well beyond the body's actual need of them for health and maintenance. John explains:

\begin{quotex}
The freedom which allows the man to turn inwards, without hindrance from sensible images, as often as he wills and thinks upon his God. This means that a man must be indifferent to gladness and grief, profit and loss, rising and falling, to strange cares, to delight and to dread, and never be attached to any creature. 

\end{quotex}
These three factors

\begin{itemize}
\item grace 
\item unity of the spirit 
\item freedom from sensible images 
\end{itemize}
are the 

\begin{quotex}
foundation and the beginning of the inward practice and the inward life. 

\end{quotex}
\paragraph{The Three Comings of Christ}
As an example of a symbolic teaching that is usually understood literally, we can look at John's teaching on the three comings of Christ:

\begin{quotex}
These are the three comings of Christ, in inward exercises. We will now explain and set forth each coming separately. Attend therefore with diligence; for he who never has himself felt or experienced this shall not easily understand it. 

\end{quotex}
John returns often to this topic, in increasing detail. For our purposes, we can focus on the first teachings. But pay careful attention to his words. One cannot reach an understanding through an intellectual argument, nor is it just a matter of blind faith. Rather, the teachings on the three comings of Christ needs to be experienced directly in order to understand them.

\textbf{First Coming}. The first coming is the incarnation, or the ``historical" Jesus. We can skip past this, since its meaning is clear. This usually deteriorates into inconclusive debates about the ``real" Jesus. Yet, paradoxically, this really can't be fully understood until the second and third comings are understood.

\textbf{Second Coming}. The second coming takes place ``in the heart", that is, it begins with second birth of the Logos in the soul. We have addressed this many times previously. Yet John makes clear that the second coming is not a one-time event. Rather, it takes place daily bringing ``new graces and new gifts". These graces make a man ``just", so he can see himself exactly as he is.

\textbf{Third Coming}. The third coming happens at the Judgment at the hour of death. Here, justice is forced on the soul if it hadn't been previously prepared by the second coming. In that case, the judgment will appear harsh, as though coming from the outside. John explains: ``the soul must then account for every word spoken and for every deed done, before the Eternal Truth."

\paragraph{Exoteric Postscript}
As should be clear by now, the religion described by John of Ruysbroeck is not intended for ``children", as \textbf{Joris-Karl Huysmans} described it. Rather, it is possible only for the spiritually mature, those who have turned away from childish things, who have tasted fully the bitterness of life and seek to transcend it.

On one side, there is the exoteric battle between the defenders and the opponents of religion. They both presume the same frame of reference. For them, there is the quest for the perfect argument that will be the ``defeater" of the opponent's position. The wisdom of the world for them is decisive: perhaps a new scientific discovery, technical advance, biblical exegesis, etc., is the answer. Or else, there is the quest for signs and wonders, or which side will manifest the will to power, whether of God or of man.

The esoteric battle, on the other hand, is an interior affair, or spiritual combat. It is not visible to the world. A fortiori, in today's world, a man can be engaged in this battle daily yet appear totally ordinary to the eyes of the world. Now, there is a pseudo-tradition abounding that rejects the True Religion of the Western Tradition in favor of some imaginary past. Keep in mind that Tradition is not a return to the past, but rather a return to what has always been. The pseudo-traditionalists therefore themselves need to participate in the spiritual combat. There is no difference, they must also dominate the animal soul through the powers of the intellect. Their path is all the more difficult without the guides or graces that make it possible.



\flrightit{Posted on 2016-01-14 by Cologero }

\begin{center}* * *\end{center}

\begin{footnotesize}\begin{sffamily}



\texttt{Matt A on 2016-01-15 at 09:42 said: }

You quote John as saying that man must ``never be attached to any creature." This got me thinking, because it made me realize that I have a lot of work to do here. I asked myself, should we abandon (literally or even just emotionally) our spouses and children due to our attachments towards them? I think not. What attachments and karmic knots we have created should be untied with love (as Mouravieff suggests), so that the relationship can continue (or not) through the guidance of love and the will of God, instead of through attachment. We must be on guard that our emotional needs towards others don't supersede God's will. Including the desire to please and comfort our closest of family, and even ourselves.


\hfill

\texttt{aegishjalmer on 2016-01-16 at 12:04 said: }

A wonderful post! Could you please elaborate on the pseudo tradition you mention, at least it's outlines?

Thank you


\hfill

\texttt{Cologero on 2016-01-17 at 10:38 said: }

@aegishjalmer Some clues:

\begin{itemize}
\item The most common indication are those who reject tout court the bearer of Tradition in the West, viz., the Medieval tradition. Granted, this is not the same as the admission that many of its traditional elements have been forgotten. 
\item They typically expect a mass intellectual conversion based, of course, on their own ideals. Rather, a movement has to begin with a smaller group. 
\item They consciously or unconsciously adopt many of the attitudes of the modern world, while claiming they are somehow ``traditional". Tradition is not a return to some part of the past, but rather to that which has eternal value. 
\item Their organizations are not hierarchically arranged on traditional principles. The pretenders to leadership positions would have been marginalized in a real traditional society. 
\item Grown men pose for publicity stills while holding a cat. 
\end{itemize}

\hfill

\texttt{William Zeitler on 2018-01-15 at 11:36 said: }

Can you suggest a translation of \textit{The Adornment of Spiritual Marriage}? Your translation sounds wonderful, but may not be available. Thanks!


\hfill

\texttt{Cologero on 2018-01-16 at 08:33 said: }

I use the translation by C. A. Wynschenk, but I am not familiar with any others. My ``translation" is more interpretive that a strict translation. I'm trying to make the point that, if you understand the Hermetic teachings, you will ``read" the literal text in a different way.


\hfill

\texttt{Petrus on 2018-01-16 at 13:26 said: }

While the Wynschenk translation is freely available online, there is a modern version whose employment of English is less stilted than the earlier one, and is an edition in the Classics of Western Spirituality series: ``John Ruusbroec: The Spiritual Espousals, The Sparkling Stones, and Other Works."


\end{sffamily}\end{footnotesize}
